% Chapter 5 — Experimental Setup

\chapter{Experimental Setup}
\label{Chapter5}

This chapter describes the experimental design for the multi-model comparative benchmark study. It covers the dataset specification, hardware infrastructure, evaluation metrics, and the protocols used to ensure reproducibility and fairness across all model comparisons.

%----------------------------------------------------------------------------------------

\section{Dataset}

\subsection{Training, Validation, and Test Splits}

All models were trained on the same \textbf{training partition} and evaluated on the same \textbf{held-out test partition} to enable direct comparison. No model had access to test data during training or hyperparameter selection. The dataset details are:

\begin{itemize}
  \item \textbf{Training Set}: 15,786 cases, 4,933 MIND rubrics, 80,049 annotated entities.
  \item \textbf{Validation Set} (early stopping): 1,973 cases, 1,973 rubrics, 9,850 entities.
  \item \textbf{Test Set} (benchmark): 1,973 cases, 1,973 rubrics, 9,955 entities.
\end{itemize}

The 80/10/10 split uses \textbf{stratified deficit-balancing}: each rubric is guaranteed to appear in the validation and test sets, preventing zero-shot evaluation on unseen rubric categories.

\subsection{Label Distribution in Test Set}

\begin{table}[h]
\caption[Test Set Label Distribution]{7-Dimension Label Distribution in the Held-Out Test Set.}
\label{tab:labeldist}
\centering
\begin{tabular}{llrr}
\toprule
\tabhead{Label} & \tabhead{Dimension} & \tabhead{Entity Count} & \tabhead{Percentage} \\
\midrule
MENT & Mental/Emotional & 4,630 & 46.5\% \\
SEN & Sensation & 1,519 & 15.3\% \\
LOC & Location & 1,272 & 12.8\% \\
MOD\_AGG & Aggravation & 803 & 8.1\% \\
TEMP & Temporal & 941 & 9.5\% \\
CONC & Concomitant & 686 & 6.9\% \\
MOD\_AMEL & Amelioration & 104 & 1.0\% \\
\midrule
\textbf{Total} & & \textbf{9,955} & 100\% \\
\bottomrule
\end{tabular}
\end{table}

The class imbalance (MENT 46.5\% vs MOD\_AMEL 1.0\%) reflects the clinical distribution of MIND chapter symptoms, where mental/emotional entities dominate.

%----------------------------------------------------------------------------------------

\section{Hardware Infrastructure}

All model training and evaluation experiments were conducted on the following hardware, accessed via IIIT Allahabad's research computing infrastructure:

\begin{table}[h]
\caption[Hardware Specification]{Experimental Hardware Specification.}
\label{tab:hardware}
\centering
\begin{tabular}{ll}
\toprule
\tabhead{Component} & \tabhead{Specification} \\
\midrule
GPU & NVIDIA RTX 5000 Ada Generation \\
VRAM & 32 GB GDDR6 ECC \\
CUDA Version & 12.4 \\
PyTorch Version & 2.6.0 \\
CPU & Intel Xeon (64-core) \\
RAM & 128 GB DDR5 ECC \\
Storage & 2 TB NVMe SSD \\
\bottomrule
\end{tabular}
\end{table}

Total GPU runtime for all experiments (synthetic data generation + 4 encoder fine-tunings + 2 LLM QLoRA fine-tunings): approximately 48 hours.

%----------------------------------------------------------------------------------------

\section{Evaluation Metrics}

\subsection{Token-Level BIO F1 (Primary NER Metric)}

Token-Level BIO F1 is computed using the \textbf{seqeval} library \cite{seqeval2018} with the strict IoB2 evaluation protocol:
\begin{itemize}
  \item A predicted entity span is correct if and only if its start position, end position, and label all exactly match the ground truth annotation.
  \item Partial overlaps are penalized as false positives and false negatives.
\end{itemize}

Precision, Recall, and F1 are computed macro-averaged across all 7 dimension labels:

\begin{equation}
F1 = \frac{2 \times \text{Precision} \times \text{Recall}}{\text{Precision} + \text{Recall}}
\label{eq:f1}
\end{equation}

\subsection{ChromaDB Retrieval Metrics}

Dense retrieval performance is measured at multiple cutoffs:
\begin{itemize}
  \item \textbf{Recall@K}: The fraction of test cases where the ground-truth rubric appears in the top-$K$ retrieved results, for $K \in \{1, 5, 10, 20\}$.
  \item \textbf{Mean Reciprocal Rank (MRR)}: The mean of the reciprocal rank of the first relevant rubric across all test queries.
\end{itemize}

\begin{equation}
\text{MRR} = \frac{1}{|Q|} \sum_{q=1}^{|Q|} \frac{1}{\text{rank}_q}
\label{eq:mrr}
\end{equation}

\subsection{Downstream Repertorization Accuracy}

Simillimum Remedy Accuracy measures whether the ground-truth remedy (the rubric's primary indicated remedy) appears in the Top-$K$ ranked candidates after the full pipeline (NER + retrieval + ranking):
\begin{itemize}
  \item \textbf{Simillimum Accuracy@1}: Top-1 Accuracy.
  \item \textbf{Simillimum Accuracy@3}: Top-3 Accuracy.
  \item \textbf{Simillimum Accuracy@5}: Top-5 Accuracy.
\end{itemize}

\subsection{LLM-Specific Metrics (BioMistral / Qwen)}

For the QLoRA-fine-tuned LLMs, two additional metrics characterize output quality:
\begin{itemize}
  \item \textbf{JSON Schema Validity Rate}: Percentage of model outputs that parse as valid JSON conforming to the \code{SymptomProfile} schema.
  \item \textbf{Clinical Grounding Rate}: Percentage of extracted entity strings that exist as verbatim substrings of the input patient narrative (hallucination detection).
\end{itemize}

%----------------------------------------------------------------------------------------

\section{Experimental Protocols}

\subsection{Reproducibility Controls}

To ensure reproducibility across all experiments:
\begin{itemize}
  \item Random seeds: \code{np.random.seed(42)}, \code{torch.manual\_seed(42)}, \code{random.seed(42)}.
  \item Deterministic CUDA operations: \code{torch.backends.cudnn.deterministic = True}.
  \item All hyperparameters frozen in \code{configs/model.yaml} and \code{configs/generation.yaml}.
  \item Identical tokenizer (AutoTokenizer from each model's official checkpoint) for all experiments.
\end{itemize}

\subsection{Evaluation Protocol for ChromaDB Retrieval}

For retrieval evaluation, extracted entity spans from each NER model are used to formulate ChromaDB query strings:
\begin{enumerate}
  \item Each extracted entity span is embedded using \texttt{all-MiniLM-L6-v2}.
  \item Top-20 rubrics retrieved per entity span.
  \item Retrieved rubric sets pooled across all entities from the same case.
  \item Ground-truth rubric (the rubric used to generate the synthetic case) checked against pooled retrieved set at each cutoff $K$.
\end{enumerate}

\textbf{DeBERTa-v3 Special Protocol}: DeBERTa-v3's BIO F1 on the test set was 0.00\% (failure to converge on the BIO labeling task due to domain gap). Its retrieval metrics were therefore evaluated by directly embedding the full patient narrative as a single query string against ChromaDB, bypassing the NER stage --- constituting a ``dense narrative retrieval'' mode that bypasses entity extraction.

\subsection{Statistical Reporting}

All reported metrics are computed on the full held-out test set ($N = 1,973$ cases) without cross-validation averaging, as the test set is sufficiently large for stable metric estimates. No significance testing is applied, as the benchmark is designed for comparative reporting rather than hypothesis testing.
